\section{Arquitectura del Sistema Híbrido}
\label{sec:cap4_arquitectura}

El reactor de bioconversión se modela como un sistema dinámico cuyas entradas ambientales (forzantes) en el instante $t$ se agrupan en el vector
\begin{equation}
\mathbf{u}(t) = \bigl[\,T(t),\;\mathrm{HR}(t),\;t\,\bigr]^{\top},
\label{eq:forzantes}
\end{equation}
donde $T$ es la temperatura del sustrato (\si{\celsius}), $\mathrm{HR}$ la humedad relativa del aire de cabeza libre (\si{\percent}) y $t$ el tiempo transcurrido desde la siembra (días). El estado biológico no observable es la biomasa viva total $B(t)$ (\si{\gram}), que el DEB predice a través de los compartimentos de carbono definidos en el Capítulo~\ref{chap:modelo}.

La arquitectura consta de seis bloques funcionales interconectados en lazo cerrado de calibración discreta (Figura~\ref{fig:arquitectura_hibrida}):

\begin{itemize}
\item[\textbf{B1}] \textbf{DEB Forward.} Simulador mecanicista que, dado $\mathbf{u}(t)$ y un vector de parámetros bioenergéticos $\boldsymbol{\theta}$, genera la tasa respiratoria teórica de dióxido de carbono $J_{\mathrm{CO_2}}(t;\boldsymbol{\theta})$ y la biomasa implícita $B(t;\boldsymbol{\theta})$.
\item[\textbf{B2}] \textbf{Sensor NDIR.} Adquiere la concentración de \ch{CO2} en el headspace del reactor, $y_{\mathrm{CO_2}}(t)$, con resolución limitada por rango (máximo \SI{6000}{ppm}) y ruido instrumental.
\item[\textbf{B3}] \textbf{Extractor de pendientes.} Procesa $y_{\mathrm{CO_2}}(t)$ en ventanas cortas post-cierre del reactor para producir el observable robusto $s_k$ (\si{ppm\per\minute}), proporcional a $J_{\mathrm{CO_2}}$.
\item[\textbf{B4}] \textbf{Comparador.} Computa el residuo de \ch{CO2} $\varepsilon_k = s_k - J_k^{\mathrm{DEB}}(\boldsymbol{\theta}_0)$, donde $\boldsymbol{\theta}_0$ es el vector de parámetros de referencia (literatura o calibración Cap~3).
\item[\textbf{B5}] \textbf{MLP Calibrador.} Red neuronal $\mathcal{N}_{\boldsymbol{\phi}}$ que, ante la discrepancia $\varepsilon_k$ y las forzantes $\mathbf{u}_k$, emite una corrección de escala $\boldsymbol{\alpha}_k$ acotada sobre el subconjunto de primera sensibilidad.
\item[\textbf{B6}] \textbf{Estimador de estado.} Integra $\boldsymbol{\theta}_k = \boldsymbol{\theta}^{\mathrm{Cap3}} \odot (1+\boldsymbol{\alpha}_k)$ en el DEB, re-simula la ventana y devuelve la biomasa estimada $\hat{B}(t_k) = B(t_k;\boldsymbol{\theta}_k)$.
\end{itemize}

\begin{figure}[htbp]
\centering
% Diagrama conceptual de bloques (placeholder TikZ)
% \shorthandoff{>} % Descomentar si se usa tikzpicture con babel spanish
\fbox{\parbox{0.85\textwidth}{\centering
\small
\textit{Placeholder: Diagrama de bloques DEB--MLP.}\\
Flujo: Sensor NDIR $\rightarrow$ Extractor $s_k$ $\rightarrow$ Comparador $\varepsilon_k$ $\rightarrow$ MLP $\boldsymbol{\alpha}_k$ $\rightarrow$ DEB Forward $\hat{B}(t_k)$.\\
Línea punteada: entrenamiento offline con pares $(\mathbf{x}_k, \varepsilon_k)$.
}}
\caption{Arquitectura híbrida DEB--MLP. El lazo de calibración opera en ventanas discretas $k$ indexadas por eventos de cierre del reactor. Durante el entrenamiento offline, el comparador alimenta al MLP con pares $(\varepsilon_k, \mathbf{u}_k)$ para ajustar $\boldsymbol{\phi}$.}
\label{fig:arquitectura_hibrida}
\end{figure}

\subsection{Flujo offline versus online}

\textbf{Etapa offline (entrenamiento).} Se ejecuta al finalizar cada ciclo experimental de 15 días. Se extraen todas las pendientes válidas $\{s_k\}_{k=1}^{K}$ de las ventanas cerradas, se simula el DEB con parámetros fijos $\boldsymbol{\theta}^{\mathrm{Cap3}}$ para obtener $\{J_k^{\mathrm{DEB}}\}_{k=1}^{K}$, se computan los residuos $\{\varepsilon_k\}_{k=1}^{K}$ y se forma el conjunto de entrenamiento $\mathcal{D}_{\mathrm{train}} = \{(\mathbf{x}_k, \varepsilon_k)\}_{k=1}^{K}$, con $\mathbf{x}_k = [\varepsilon_k, T_k, \mathrm{HR}_k, t_k]^{\top}$. El MLP se entrena minimizando la función de pérdida compuesta definida en la Sección~\ref{sec:cap4_calibrador}.

\textbf{Etapa online (inferencia).} Por cada nueva ventana de muestreo durante el ciclo activo:
\begin{enumerate}
\item Medir $s_k$ del sensor NDIR y registrar $\mathbf{u}_k$.
\item Simular el DEB con $\boldsymbol{\theta}^{\mathrm{Cap3}}$ para obtener $J_k^{\mathrm{DEB}}$.
\item Calcular $\varepsilon_k = s_k - J_k^{\mathrm{DEB}}$.
\item Formar $\mathbf{x}_k$ y evaluar el MLP: $\boldsymbol{\alpha}_k = \mathcal{N}_{\boldsymbol{\phi}}(\mathbf{x}_k)$.
\item Actualizar parámetros: $\boldsymbol{\theta}_k = \boldsymbol{\theta}^{\mathrm{Cap3}} \odot (1+\boldsymbol{\alpha}_k)$.
\item Re-simular el DEB con $\boldsymbol{\theta}_k$ en la ventana corta y emitir $\hat{B}(t_k)$.
\end{enumerate}

La latencia de esta secuencia, ejecutada en el microcontrolador ESP32 o en un backend ligero, es inferior a \SI{5}{\second} por ventana, cumpliendo el requisito de utilidad operativa establecido en el Capítulo~\ref{chap:metricas}.
